\textcolor{black}{\section{Introducción}}

\noindent El transporte de electrones en gases es un fenómeno fundamental en la física de plasmas y de descargas eléctricas, ya que determina la forma en que las partículas cargadas adquieren energía del campo eléctrico y la transfieren al medio mediante colisiones. La descripción de este transporte depende de la probabilidad de interacción entre los electrones y las moléculas o átomos del gas, la cual se cuantifica mediante la sección eficaz (\textit{cross section}) de dispersión. \cite{Urquijo}

Para estudiar plasmas de baja temperatura, se usan las técnicas de enjambre, en estas técnicas no nos importa lo que le sucede a un solo electrón, sino lo que ocurre en el fenómeno colectivo del enjambre de electrones, en particular estaremos usando la técnica pulsada de Townsend, esta técnica se basa en dos principios fundamentales. El efecto fotoeléctrico y el efecto Ramsauer--Townsend.

Uno de los fenómenos más importantes en la dispersión de electrones de baja energía es el efecto Ramsauer-Townsend. Este efecto consiste en la aparición de un mínimo en la sección eficaz de dispersión elástica de algunos gases nobles para un intervalo específico de energías electrónicas. Como consecuencia, los electrones experimentan una menor probabilidad de colisionar con los átomos del gas, aumentando su camino libre medio y modificando significativamente sus propiedades de transporte. \cite{Olmo}

Experimentalmente, estos fenómenos pueden estudiarse mediante la técnica de Townsend pulsado. En este método, un pulso láser ultravioleta incide sobre el cátodo de un capacitor metálico produciendo, por efecto fotoeléctrico, un pulso inicial de electrones. Bajo la acción de un campo eléctrico uniforme, estos electrones forman un enjambre, debido a la ionización de las partículas, que se desplaza a través de la mezcla gaseosa , generando una corriente transitoria cuya forma depende de los procesos de deriva, difusión e ionización que experimentan durante su trayectoria. \cite{Olmo}

El efecto Ramsauer-Townsend puede comprobarse experimentalmente, observando y midiendo la dependencia que hay entre la velocidad de deriva de los electrones y el aumento en el campo eléctrico producido por el capacitor y con las diferentes presiones, propiciando así las colisiones y la ionización del gas. Esto también nos permite observar el punto mínimo de la sección eficaz que describe este efecto. \cite{Swarm}
\noindent 